% concatenated from main.tex
\documentclass[reqno, a4paper]{amsart}

\usepackage{graphicx} % Required for inserting images
% Make references clickable
\usepackage[hidelinks]{hyperref}
% 
\hypersetup{colorlinks=true,linkcolor=blue,citecolor=blue}
\usepackage[utf8]{inputenc}
\usepackage[german, english]{babel}
\usepackage[T2A,OT2,OT1]{fontenc}

\newcommand{\lc}{\preceq_{\textnormal{c}}}

\newcommand{\ld}{\preceq_{E}}


\newcommand{\lbd}{\preceq_{\bar D}}

\newcommand{\lic}{\preceq_{icx}}

\usepackage{datetime}

\newdateformat{monthyeardate}{%
  \monthname[\THEMONTH], \THEYEAR}
\usepackage{amsfonts}

\usepackage{bbm}
\usepackage{amssymb}
%
%\usepackage[english, russian]{babel}
\usepackage{amsthm}
\usepackage[dvipsnames]{xcolor}
\newtheorem{theorem}{Theorem}[section]
\newtheorem{corollary}[theorem]{Corollary}
\newtheorem{lemma}[theorem]{Lemma}
\usepackage{mathtools}
\usepackage{accents}
\usepackage{enumitem}
\newcommand{\dbtilde}[1]{\accentset{\approx}{#1}}
\newtheorem{thmm}{Theorem}
%\newcommand\myeq{\mathrel{\stackrel{\makebox[0pt]{\mbox{\normalfont\tiny Thm 1}}}{=}}}
\newcommand{\conv}{\mathrm{conv}\,}
\newcommand{\convr}{\mathrm{conv}_R\,}

\newcommand{\iconv}{\mathrm{iconv}\,}
\newcommand{\PP}{\mathcal{P}}
\newcommand{\WW}{\mathcal{W}}
\newcommand{\Pt}{\mathcal{P}_{p}(\mathbb{R}^d)}
\newcommand{\Pcs}{\mathcal{P}_{\text{cs}}(\mathbb{R}^d)}
\newcommand{\cP}{\mathcal{P}}
\newcommand{\Pc}{\cP}
\renewcommand{\P}{\cP}
\newcommand{\env}{\mathrm{env}_E\,}

\newcommand{\midauto}{\,\middle\vert\,}
\newcommand{\R}{\mathbb R}
\theoremstyle{definition}
\newtheorem{definition}[theorem]{Definition}
  \theoremstyle{definition}
\newtheorem{example}[theorem]{Example}
\newtheorem{prop}[theorem]{Proposition}
\theoremstyle{definition}
\newtheorem{remark}[theorem]{Remark}


\newcommand{\cpl}{\Pi}
\newcommand{\cplm}{\cpl_\mathrm{M}}

\newcommand{\LSC}{\mathcal{LSC}}
\newcommand{\mean}{\mathrm{mean}}


\title{Weak Optimal Transport: When is the Dual Potential convex?}
\author{Filip Pramenković}
\date{\monthyeardate\today}

\begin{document}

\begin{abstract}
 Weak optimal transport generalizes the classical theory of optimal transportation to nonlinear cost functions and covers a range of problems that lie beyond the traditional theory --- including entropic transport, martingale transport, and applications in mechanism design. As in the classical case, the weak transport problem can also be written as a dual maximization problem over a pair of conjugate potentials. 

 We identify sharp monotonicity conditions on the cost under which the dual problem can be restricted to \emph{convex} potentials. This framework unifies several known results from the literature, including barycentric transport, martingale Benamou--Brenier, the multiple-good monopolist problem, Strassen's theorem, stochastic order projections and the classical Brenier theorem. %Moreover, we isolate the essential properties needed for the dual simplification and discuss a general setting beyond convex functions.
 
 
 
%  This allows to short, unified several known results from the literature, including martingale Benamou--Brenier, barycentric transport, the multiple-good monopolist problem, the proof of Strassen's theorem and stochastic order projections. Lastly, we isolate the essential properties needed for the dual simplification and discuss a general framework beyond convex functions 


% %The weak transport problem arose from questions in geometric inequalities and was formally introduced in 2014 by Gozlan, Roberto, Samson, and Tetali. It allows for nonlinear cost structures and covers a range of problems that lie beyond the traditional theory --- including entropic transport, martingale transport, and applications in mechanism design.





%     Weak optimal transport generalizes the classical theory of optimal transportation to nonlinear cost functions, encompassing variants such as martingale and entropic optimal transport. Since weak transport problems amount to convex optimization tasks, they naturally satisfy a duality principle. The associated dual problem is an optimization task over a class of functions, prompting us to ask how much this function space can be refined.
    
%     We identify equivalent monotonicity conditions under which the dual problem can be restricted to convex and increasing convex functions on $d$-dimensional Euclidean space. We additionally show attainment of the dual problem under certain regularity assumptions. 
    
%     The developed framework unifies several known results from the literature, including martingale Benamou--Brenier, barycentric transport, the multiple-good monopolist problem, the proof of Strassen's theorem and stochastic order projections. Lastly, we isolate the essential properties needed for the dual simplification and discuss a general framework beyond convex functions and the Euclidean setting.
\end{abstract}

\maketitle


\section{Introduction}
%The field of weak optimal transport (WOT), introduced by \text{Gozlan, Roberto, Samson} and \text{Tentali} \cite{gozlan2015kantorovich}, is an extension of classical transport theory to non-linear cost functions. It encompasses important applications that fall outside the classical transport framework, while still allowing for the same key existence and duality results. 

Weak optimal transport was introduced by Gozlan, Roberto, Samson and Tetali \cite{GoRoSaTe14} to extend classical transport theory to nonlinear cost functions. It  encompasses important applications that fall outside the classical transport framework while still allowing for the same basic existence and duality results, see e.g.\  \cite{GoRoSaSh18, AlBoCh18, BaPa19, BePaRiSc25}.

To introduce the weak transport problem, we set up some notation.
We consider Polish metric spaces $(X,d_X), (Y, d_Y)$ equipped with Borel probability measures $\mu, \nu$, respectively.
%Given a Polish metric space $Z$, we denote by $\cP(Z)$ the set of probabilities on $Z$, equipped with weak convergence, and by $\cP_p(Z)$, $p\in [1,\infty)$ probabilities with finite $p$-moment, equipped with $p$-weak convergence. 
% Throughout we consider Polish probability spaces $(X, \mu)$, $(Y,\nu)$. 
We write $\cpl(\mu, \nu)$ for the set of all \emph{couplings} or \emph{transport plans}, that is, probabilities on \(X \times Y\) which have \(X\)-marginal \(\mu\) and \(Y\)-marginal \(\nu\). For $p\geq 1$ we write $\Pc_p(Y)$ for the set of probabilities with finite $p$-moment, equipped with $p$-weak convergence. 
 %the purpose of defining moment assumptions we fix a complete compatible metric $d$ on $Y$ and $p\in [1,\infty)$. 
Given a lower semicontinuous (lsc) cost function $C:X\times \Pc_p(Y)\to [0,\infty]$ that is convex in the second argument, the primal weak transport problem is given by
\begin{align}\label{eq:PrimalWOT}
    \inf_{\pi\in\Pi(\mu,\nu)}\int C(x,\pi_x)d\mu(x),
\end{align}
where $\pi_x$ denotes the disintegration of $\pi$ with respect to its first marginal. As in classical optimal transport, lower semicontinuity of $C$ implies that \eqref{eq:PrimalWOT} is attained and equals the dual problem
\begin{align}
    \sup_{\psi \in {\mathcal{C}}_{b,p}(Y)} \mu(\psi^C) - \nu(\psi), \quad \psi^C(x):= \inf_{\rho\in \Pc_p(Y)}\rho(\psi) + C(x,\rho). 
\end{align}
We write $\mathcal C_{b,p}(Y)$ for the subset of continuous functions bounded from below by a constant and bounded from above by a multiple of $1+d_Y(\cdot,y_0)^p$ for some $y_0\in Y$.  

In many applications of WOT-duality, the dual problem simplifies significantly. For example, \ given $\mu, \nu\in \Pc_1(\R^d)$,  we have 
\begin{align}
    \inf_{\pi\in \cpl(\mu, \nu)} \int \chi_{x= \mean(\pi_x)} \, d\mu(x) = \sup_{ \substack{\psi \textrm{ convex} \\ \psi \in \mathcal C_{b,1}(\R^d)} } \mu(\psi)- \nu(\psi),
\end{align}
where $\chi_{x= \mean(\pi_x)}=0$, if $x= \mean(\pi_x)$, and $\infty$ otherwise. This duality result was first noticed in \cite{GoRoSaTe14} and directly implies Strassen's theorem on the existence of martingales between marginals in convex order, see Theorem \ref{strass} below for details.

To further illustrate the point, for $\mu, \nu \in \Pc_1(\R)$ we have 
\begin{align}
    \inf_{\pi\in \cpl(\mu, \nu)} \int (x-\mean(\pi_x))_+ \, d\mu(x) = \sup_{\substack{\psi \textrm{ increasing convex}\\ 1\textrm{-Lipschitz}}} \mu(\psi)- \nu(\psi).
\end{align}
This is referred to as the increasing convex Kantorovich--Rubinstein formula in \cite{BePaRiSc25}. 

We highlight that in both examples it suffices to optimize  over a smaller class of functions in the dual problem, namely convex and increasing convex functions. In fact, the same is true for several further problems in the weak transport framework: the original convex Kantorovich--Rubinstein formula \cite{GoRoSaTe14}, the Brenier--Strassen theorem \cite{GoJu18}, the dual formulation of the martingale Benamou--Brenier problem \cite{backhoffveraguas2025existencebassmartingalesmartingale}, optimal mechanism design \cite{DaDeTz17}, and the Gangbo--McCann--Strassen theorem \cite{BePaRiSc25}. 

The purpose of this note is to provide a unified explanation of the aforementioned dual problem simplifications. We specialize to the case of $Y=\R^d$, which coincides with the listed examples and simplifies our discussion. It turns out that \textit{stochastic orders} are the correct tool for explaining this phenomenon. Recall that $\rho_1, \rho_2 \in\PP_p(\R^d)$ are in  \textit{convex order}, in signs $\rho_1\lc \rho_2$, if $\rho_1(f) \leq \rho_2(f)$ for all convex functions $f:\R^d\to\R$. 

Our main result relates convexity of the dual potentials in weak transport problems to a monotonicity condition for the cost function:
\begin{theorem} \label{mainthm}
    Let $\mu\in\PP(X)$, $\nu \in \PP_p(\R^d)$ and $C:X\times\PP_p(\R^d)\to[0,\infty]$ be lsc and convex in the second argument. If $C$ is $\lc$-decreasing in the second argument, then 
\begin{align} \label{mainthmdispl}
		\inf_{\pi\in\Pi(\mu,\nu)} \int C(x,\pi_x)d\mu(x)= \sup_{\substack{\psi \in  {\mathcal{C}}^{}_{b,p}(\R^d) \\ \psi\textnormal{ convex}}} \mu(\psi^C) - \nu(\psi).
	\end{align}
    Conversely, if (\ref{mainthmdispl}) holds for all $\mu\in\PP(X)$, $\nu\in\PP_p(\R^d)$, then $C$ is $\lc$-decreasing in the second argument.
\end{theorem}

Furthermore, in Theorem \ref{thmatt} we show that under suitable regularity assumptions, the supremum in (\ref{mainthmdispl}) is attained by a $\nu$-integrable convex function $\psi^{\text{opt}}:\R^d\to(-\infty,\infty]$. More precisely, to guarantee attainment we require that the cost function does not grow too rapidly and that it exhibits some mild continuity in the second argument, see Section \ref{sec att}.

Analogous duality and attainment results hold for coordinatewise increasing convex functions and the associated \textit{increasing convex order} denoted by $\lic$, see Theorem \ref{mainthmic} and Theorem \ref{thmatt2}. 

In fact, the simplification of the dual problem can be extended to general Polish spaces and orders defined by \textit{stable cones}, an abstract class of functions that captures the relevant properties of convex and increasing convex functions, see Section \ref{nesto}.







\medskip

% \begin{remark}
%     In the case of a compact Polish space $Y$, this principle can be extended to any convex cone $D\subseteq {\mathcal{C}}(Y)$ that contains all the constant functions and is closed under taking finite maxima. Indeed, $C$ is $\ld$-decreasing in the second argument, if and only if, the dual can be restricted to functions in $D$. See Section \ref{CompactCase}. 
% \end{remark}


% \section{Introduction}

% The classical theory of optimal transport has its origins in the works of \text{Monge} \cite{monge1781mémoire} and \text{Kantorovich} \cite{Ka42}. Let $\mu$ and $\nu$ be probability measures on Polish spaces $X$ and $Y$, respectively. Given a measurable cost function $c:X\times Y\to(-\infty,\infty]$, the Kantorovich problem and its dual are given by
% \begin{align} \tag{OT} \label{prva}
%     \inf_{\pi\in\Pi(\mu,\nu)}\int c(x,y)d\pi(x,y)=\sup_{\substack{\phi\in L^1(\nu)}} \nu(\phi)-\mu( \phi^c ),
% \end{align}
% where $\Pi(\mu,\nu)$ denotes the set of all couplings between $\mu$ and $\nu$, that is, probability measures on the product space $X\times Y$ with first marginal distribution $\mu$ and second marginal distribution $\nu$. The duality in (\ref{prva}) holds under mild assumptions (see, e.g., \cite[Theorem 5.10]{villani2008optimal}). The function $\phi^c$ denotes the $c$-conjugate of $\phi$, defined as $\phi^c(x):=\sup_{y\in Y}[\phi(y)-c(x,y)]$. Note that we employ the notational shorthand $\nu(\phi):=\int\phi d\nu$ for integration throughout the paper. 

% Weak optimal transport, introduced in \cite{gozlan2015kantorovich} by  \text{Gozlan}, \text{Roberto}, \text{Samson} and \text{Tetali}, is a generalization of the classical theory of transportation. Let $X,Y$ be Polish spaces and fix a compatible metric $d$ on $Y$. For the purpose of defining moment assumptions, fix $p\in[1,\infty)$. We write $\PP_p(Y)$ for the space of probability measures on $Y$ with finite $p$-th moment. 

% The weak optimal transport problem associated to a cost function $C:X\times\PP_p(Y)\to(-\infty,\infty]$ and its dual problem are given by
% \begin{equation}
% \label{WOT}
% \tag{WOT}
%     \inf_{\pi\in\Pi(\mu,\nu)}\int C(x,\pi_x)d\mu(x)=\sup_{\psi \in {\mathcal{C}}_{b,p}(Y)} \mu(\psi^C) - \nu(\psi),
% \end{equation}
% where $\pi_x$ denotes the disintegration of the coupling $\pi\in\Pi(\mu,\nu)$ such that $d\pi(x,y)=d\pi_x(y)d\mu(x)$. The duality principle for weak optimal transport holds under mild assumptions (see  \cite{veraguas2019existence} or Theorem \ref{th:kantorovichduality}). We write $\psi^C$ for the $C$-conjugate of $\psi$ defined by $\psi^C(x):=\inf_{\rho\in\PP_p(Y)} \rho(\psi)+C(x,\rho)$ and ${\mathcal{C}}_{b,p}(Y)$ for the set of continuous function on $Y$ that are bounded from below and of at most $p$-growth.

% The main contribution of this work lies in identifying a necessary and sufficient condition for restricting the dual problem to a class of functions associated with a given convex cone of continuous functions. We introduce this principle by exploring two prominent examples.
% \subsection{Restriction to convex functions: Martingale Benamou-Brenier}
% Let $\mu,\nu\in\PP_2(\R^d)$ be in convex order, $\mu\lc\nu$ (that is, $\mu(f)\leq\nu(f)$ for all convex functions $f$), and let us consider the martingale Benamou-Brenier problem and its dual 
% \begin{equation}
% \tag{MBB}\label{mbb}\sup_{\substack{M_0\sim\mu, M_1\sim\nu \\ M_t=M_0+\int_0^t\sigma_sdB_s}} \mathbb E \left[\int_0^1 \text{tr}(\sigma_t)dt\right]=\inf_{\substack{\psi\in L^1(\nu)\\\psi \text{ convex}}}\nu(\psi)-\mu((\psi^**\gamma)^*),
% \end{equation}
% where $\gamma$ denotes the standard normal distribution on $\R^d$, the symbol $*$ is the convex conjugate when used in superscript, and otherwise denotes the convolution operator. 
% Under some regularity assumptions, the unique maximizer is given by the \textit{Bass martingale} 
% $$M_t:=\mathbb E \left[ \nabla\psi^*_{\text{opt}}(B_1)\mid \sigma(B_s:s\leq t)\right], \quad \quad 0\leq t\leq 1,$$
% where $B_t$ is Brownian motion with initial distribution $B_0\sim\nabla(\psi^*_{\text{opt}} *\gamma)^*(\mu)$ and $\psi_{\text{opt}}$ is a minimizer (in a relaxed sense) of the right-hand side of (\ref{mbb}) (see \cite{10.1214/20-AOP1422,backhoffveraguas2025existencebassmartingalesmartingale}).

% The restriction to convex functions in (\ref{mbb}) follows from the equivalent duality of the \textit{static version} of the martingale Benamou-Brenier problem
% \begin{equation}
% \label{SMBB}
% \tag{SMBB}
%     \sup_{\pi\in\text{MT}(\mu,\nu) } \int \text{MCov}(\pi_x,\gamma)d\mu(x)=\inf_{\substack{\psi \in C_q(\R^d) \\ \psi \text{ convex}}}\nu(\psi)-\mu(\varphi^\psi),
% \end{equation}
% where $\text{MT}(\mu,\nu)$ denotes the set of martingale couplings between $\mu$ and $\nu$,  $C_q(\R^d):={\mathcal{C}}_{2,p}(\R^d)+|\cdot|^2/2$ and $\text{MCov}(\rho_1,\rho_2):=\max_{q \in \Pi(\rho_1,\rho_2) } \int_{\R^d\times \R ^d}yz dq(y,z)$. Notice that this is a weak transport problem with the cost function $$C(x,\rho)=  \begin{cases} 
%         -\text{MCov}(\rho,\gamma), & x=\text{bary} (\rho) \\
%       \infty, & \text{otherwise}, 
%    \end{cases}
%   $$
%   and the function $\varphi^\psi$ occurring in (\ref{SMBB}) is different notation for $\psi^C$, the $C$-conjugate of $\psi$. Hence (\ref{WOT}) implies the duality in (\ref{SMBB}), up to the restriction to convex functions. 
  
%   It turns out that the ability to restrict the dual problem to convex functions is intimately tied to the monotonic structure of the cost function. More precisely, such a restriction is possible if and only if $C$ is $\lc$-decreasing in the second argument, that is $\rho_1\lc\rho_2\implies C(x,\rho_1)\geq C(x,\rho_2)$ for all $x\in X$. We capture this principle in the following theorem.

% \begin{theorem} \label{convthm}
%     Let $\mu\in\PP(X)$, $\nu \in \PP_p(\R^d)$ and let $C:X\times\PP_p(\R^d)\to(-\infty,\infty]$ be jointly lsc\footnote{lower semicontinuous}, convex in the second argument and bounded from below. If $C$ is $\lc$-decreasing in the second argument, then 
% \begin{align} \label{mainthmdispl2}
% 		\inf_{\pi\in\Pi(\mu,\nu)} \int C(x,\pi_x)d\mu(x)= \sup_{\substack{\psi \in {\mathcal{C}}^{}_{b,p}(\R^d)\\\psi\textnormal{ convex}}} \mu(\psi^C) - \nu(\psi).
% 	\end{align}
%     Conversely, if (\ref{mainthmdispl2}) holds for all $\mu\in\PP(X)$ and $\nu\in\PP_p(\R^d)$, then $C$ is $\lc$-decreasing.
% \end{theorem}

% \begin{remark}
%     The lower bound assumption can be relaxed, see Theorem \ref{thmmod1}.
% \end{remark}

% \subsection{Restriction to increasing convex functions: Multiple-good Monopolist}
% In \cite{daskalakis2017strongdualitymultiplegoodmonopolist}  \text{Daskalakis, Deckelbaum,} and\text{ Tzamos} study the economic problem of optimizing revenue for a multiple-good monopolist by phrasing it in terms of an optimal transportation problem and utilizing duality theory to characterize optimal mechanisms. Furthermore, it can be stated in terms of a weak optimal transport problem (see \cite{backhoffveraguas2020applicationsweaktransporttheory}). 

% Let $\mu,\nu\in\PP_p(\R^d)$, and fix $\theta\in{\mathcal{C}}_{b,p}(\R^d)$ be convex. We equip $\R^d$ with coordinate-wise partial order $\leq$. We have the duality
% \begin{equation}\label{mgm}\tag{MGM}
%     \inf_{\pi \in \Pi(\mu,\nu)} \int \inf_{\xi\lic\rho}\theta(x-\text{bary}(\xi))d\mu(x)=\sup_{\substack{\psi\in{\mathcal{C}}_{b,p}(\R^d) \\\psi\text{ increasing convex}} }\mu(R_\theta \psi)-\nu(\psi),
% \end{equation}
% where $\xi\lic\rho$ represents $\xi$ and $\rho$ being in \textit{increasing convex order,} i.e., $\xi(f)\leq\rho(f)$ for all functions $f$ on $\R^d$ that are (coordinate-wise) increasing and convex. Note that the weak cost function in this case is $C(x,\rho)=\inf_{\xi\lic\rho}\theta(x-\text{bary}(\xi))$, and $R_{\theta} \psi=\psi^C$.  

% Much like in the previous example, if the cost function is decreasing in the second argument with respect to $\lic$, then the dual may be restricted to increasing convex functions, and vice-versa.


% \begin{theorem} \label{iconvthm}
%     Let $\mu\in\PP(X)$, $\nu \in \PP_p(\R^d)$ and let $C:X\times\PP_p(\R^d)\to(-\infty,\infty]$ be jointly lsc, convex in the second argument and bounded from below. If $C$ is $\lic$-decreasing in the second argument, then 
% \begin{align} \label{mainthmdispl22}
% 		\inf_{\pi\in\Pi(\mu,\nu)} \int C(x,\pi_x)d\mu(x)= \sup_{\substack{\psi \in {\mathcal{C}}^{}_{b,p}(\R^d)\\\psi\textnormal{ increasing convex}}} \mu(\psi^C) - \nu(\psi).
% 	\end{align}
%     Conversely, if (\ref{mainthmdispl22}) holds for all $\mu\in\PP(X)$ and $\nu\in\PP_p(\R^d)$, then $C$ is $\lic$-decreasing.
% \end{theorem}
% \begin{remark}
%     Note that Theorem \ref{convthm} also applies in the case of (\ref{mgm}), but only gives us a restriction to convex functions. This is because any $\lic$-decreasing map is necessarily $\lc$-decreasing. 
% \end{remark}


% \bigskip

% This pattern extends beyond convex and increasing convex functions. In the main technical theorem of the paper, Theorem \ref{mainthm}, we show that the dual can be restricted to a set associated to an arbitrary convex cone of continuous functions that is closed with respect to taking maxima precisely when the cost function is decreasing with respect to the stochastic order induced by this convex cone. Specializing to the convex cone of convex functions and increasing convex functions establishes Theorem \ref{convthm} and \ref{iconvthm}, respectively.

The paper is organized as follows. Section \ref{PR} serves as a preparatory section, where we collect relevant preliminary results for our discussion. In Section \ref{MT}, we prove the central theorems on restricting the dual to (increasing) convex functions and in Section \ref{sec att} we consider under what assumptions and in what sense the dual problem is attained.  In Section \ref{APP} we apply the main theorems to recover known results from the literature and in Section \ref{nesto} we discuss a general abstract framework for the restriction problem. 










   %  \item The sets ${\mathcal{C}}^{lsc}_{p}(Y)$ and ${\mathcal{C}}^{lsc}_{b,p}(Y)$ denote the lsc analogues of ${\mathcal{C}}_{p}(Y)$ and ${\mathcal{C}}^{}_{b,p}(Y)$, respectively, where instead of continuity, we require the functions to be lsc.
    % \item We define the lower semicontinuous analogue of ${\mathcal{C}}_{p}(Y)$ as
%$${\mathcal{C}}^{lsc}_{p}(Y):=\left\{\phi:Y\to\R \mid \phi \text{ lsc, } \exists y_0\in Y, a\in\R:|\phi|\leq a(1+d(y_0,\cdot)^p)  \right\} $$ and we write
%${\mathcal{C}}^{lsc}_{b,p}(Y)$ for the set of functions in ${\mathcal{C}}^{lsc}_{p}(Y)$ which are additionally bounded from below.

%\item Let $D\subseteq {\mathcal{C}}_p(Y)$ be a convex cone, i.e., if $f,g\in D$, then $f+ g \in D$ and $\alpha f \in D$ for all $\alpha\geq0$. We additionally assume that $D$ contains all constant functions and is max-stable, that is, $f,g\in D$ implies $f\lor g \in D$. 
%     \item Define the relation $\lc$ on $\PP_p(Y)$ by
 %     \begin{equation} \label{ld def} \mu\lc\nu \text{ iff }\left(\mu(\phi)\leq\nu(\phi) \text{ for all }\phi\in  \right).\end{equation}
  %     Note that the relation $\ld$ is a preorder on $\PP_p(Y)$, that is, a reflexive and transitive binary relation. 
   %     \item For $\psi\in C(Y)$, we define its $D$-hull as
%     $$\widehat{\psi}(y):=\sup\{\phi(y) \mid \phi\in D,\phi \leq \psi \}.$$
    % \item Denote by $\bar D$ the set of all functions realized as pointwise suprema of elements of $D$. Note that $\bar D$ is a convex cone since $D$ is. Functions in $\bar D$ are in general $(-\infty,\infty]$-valued and lsc, as they are suprema of continuous functions.


\section{Preliminary Results}
\label{PR}
In this section, we lay the foundations of our main theorems by discussing and proving some auxiliary results. 
% Let us fix some notation.

% \begin{itemize}
% \item Let $X,Y$ be Polish spaces and equip $Y$ with a compatible metric $d$. Fix $p\in[1,\infty)$. 
%     \item ${\mathcal{C}}(Y)$ denotes the space of real-valued continuous functions on $Y$.
%     \item $\PP(X)$ is the space of Borel probability measures on $X$.
%      \item We write $\PP_p(Y)$ for the space of Borel probability measures $\rho$ on $Y$ such that for some (and then any) $y_0\in Y$ one has $\int d(y_0,y)^pd\rho(y)<\infty$.
%      \item The set ${\mathcal{C}}_{p}(Y)$ contains functions $\psi$ in ${\mathcal{C}}(Y)$ such that for some $y_0\in Y$ there is a constant $a\in\R$ that satisfies $|\psi(y)|\leq a (1+d(y_0,y)^p)$.
%      \item ${\mathcal{C}}_{b,p}(Y)$ denotes the set of functions $\psi$ in ${\mathcal{C}}(Y)$ such that for some $y_0\in Y$ there are constants $a,\ell\in\R$ that satisfy $\ell\leq\psi(y)\leq a (1+d(y_0,y)^p)$.
%      \item We equip the space $\R^d$ with the usual Euclidean distance, which we denote by $|\cdot|$.
%      \end{itemize}


\subsection{Convex and Increasing Convex Hulls}
 We start by recalling the definitions of the convex hull and the increasing convex hull, which play a key role in proving the main theorems.

    Let $\psi\in{\mathcal{C}}_{b,p}(\R^d)$. We denote its \textit{convex hull} as 
    $$\conv\psi(y):=\sup_{ \substack{\phi\le\psi \\ \phi \text{ convex}}} \phi(y), \;\; y\in\R^d,$$
and its \textit{increasing convex hull} as 
    $$\iconv\psi(y):=\sup_{ \substack{\phi\le\psi \\ \phi \text{ inc. convex}}} \phi(y), \;\; y\in\R^d.$$


% Simply put, the convex hull of $\psi$ is the greatest convex function majorized by $\psi$, whereas the increasing convex hull is the greatest coordinatewise increasing convex function majorized by $\psi$.

% \begin{theorem} \label{teorema gogus}{\cite[{Theorem 1.3}]{Gogus2013}} Let $T$ be a locally compact $\sigma$-compact Hausdorff space and let $D\subseteq C(T)$ be a convex cone containing all constant functions. If $\psi\in C(T)$, then  
% \begin{equation} \label{nejed}
%     \widehat{\psi}(t)=\inf\left\{\rho(\phi) \mid \delta_t\ld\rho,\; \rho\in\PP_{cs}(T)\right\},
% \end{equation}
% unless
% $\widehat{\psi}\equiv - \infty$.
    
% \end{theorem}

Recall that for $\rho_1,\rho_2\in\P_p(\R^d)$ the notation $\rho_1\lc \rho_2$ (resp. $\rho_1\lic \rho_2$), means that $\rho_1(f)\le\rho_2(f)$ for all $f:\R^d\to\R$ convex (resp. increasing convex).

The following lemma is a collection of well-known facts about the convex and increasing convex hull.


\begin{lemma} \label{repr lemm}
    Let $\psi \in {\mathcal{C}}_{b,p}(\R^d)$. Then we have 
    \begin{enumerate}[label=(\roman*)]\label{repr enac}
     \item \label{conv repr bullet}  $\conv{\psi}(y)=\inf\left\{\rho(\psi) \mid \delta_y\lc\rho,\; \rho\in\PP_p(\R^d)\right\},$
     \item \label{iconv repr bullet}  $\iconv{\psi}(y)=\inf\left\{\rho(\psi) \mid \delta_y\lic\rho,\; \rho\in\PP_p(\R^d)\right\},$
      \item \label{prop1} $\conv{\psi}\leq \psi$,
     \item \label{prop3}$\psi_1,\psi_2\in{\mathcal{C}}_{b,p}(\R^d)$ and $\psi_1\leq\psi_2$, then $\conv{\psi_1}\leq \conv{\psi_2},$
      % \item \label{prop5}if $\psi $ is convex, then $\conv{\psi}=\psi$,
  %  \item \label{prop2}$\widehat{\psi}$ is real-valued,
    \item \label{prop4}$\conv{\psi}\in{\mathcal{C}}_{b,p}(\R^d),$
    \item \label{prop6}$\conv{\psi_n}\downarrow \conv{\psi}$, as $n\to\infty$, where $\psi_n:=\psi+\frac{1}{n}|\cdot|^p.$
    \end{enumerate}
Additionally, properties \ref{prop1}--\ref{prop6} also hold for the increasing convex hull.
\end{lemma}
\begin{proof}
Since \textit{\ref{iconv repr bullet}} is less commonly cited than \textit{\ref{conv repr bullet}}, let us prove it. 
Fix $\psi\in{\mathcal{C}}_{b,p}(\R^d)$, $y\in \R^d$ and $\rho\in\P_p(\R^d)$ such that $\delta_y \lic \rho$. Then
$$\iconv\psi(y)=\sup_{ \substack{\phi\le\psi \\ \phi \text{ inc. convex}}} \phi(y)\leq \sup_{ \substack{\phi\le\psi \\ \phi \text{ inc. convex}}} \rho(\phi) \leq \rho(\psi),$$
and taking the infimum over all such $\rho$ proves the first direction. 

Now write $$f(y):=\inf\left\{\rho(\psi) \mid \delta_y\lic\rho,\; \rho\in\PP_p(\R^d)\right\}.$$
To prove the converse inequality, it suffices to show that $f$ is increasing, convex and majorized by $\psi$.

Note that choosing $\rho=\delta_y$ yields $f(y)\leq\psi(y)$. Next, we show that $f$ is convex. Fix $y_1,y_2\in\R^d$, $\alpha\in(0,1)$ and $\varepsilon>0$. Pick $\rho_1,\rho_2 \in\P_p(\R^d)$ with  $\delta_{y_1}\lic \rho_1$ and  $\delta_{y_2}\lic \rho_2$ such that $\rho_1(\psi)\leq f(y_1)+\varepsilon$ and $\rho_2(\psi)\leq f(y_2)+\varepsilon$. Then
$$f(\alpha y_1+(1-\alpha)y_2)\leq (\alpha\rho_1+(1-\alpha)\rho_2)(\psi) \leq \alpha f(y_1)+(1-\alpha)f(y_2)+ \varepsilon.$$

Lastly, we show that $f$ is increasing. Let $y_1,y_2\in\R^d$ be such that $y_1\leq y_2$ coordinatewise. In particular, we have $\delta_{y_1}\lic\delta_{y_2}$. But by definition of $f$, this immediately implies $f(y_1)\leq f(y_2).$ 

Points \textit{\ref{prop1}--\ref{prop6}} are obvious.
\end{proof}

% \begin{remark}
%     The condition $\delta_y\lc\rho$ (resp. $\delta_y\lic \rho$) occurring in  Lemma \ref{repr lemm} can equivalently be stated as $y=\mean(\rho)$ (resp. $y\leq\mean(\rho)$). Nonetheless, we prefer to use the stochastic order form to highlight the connection with the general framework outlined in  Section \ref{nesto}.
% \end{remark}

% The following lemma provides us with a few useful properties of the convex and increasing convex hull.


% \begin{proof} Points 
% \textit{\ref{prop1}} and \textit{\ref{prop3}} are obvious.

% \smallskip
% \noindent 
% \textit{\ref{prop4}} By \textit{\ref{prop1}}, $\conv{\psi}\leq\psi\leq a(1+|\cdot|^p)$. By the dual representation in Lemma \ref{repr lemm}, any constant lower bound $\ell$ of $\psi$ is clearly a lower bound of $\conv{\psi}$, so $\ell\le\conv \psi\leq a(1+|\cdot|^p)$.  The continuity of $\conv \psi$ follows from the fact that real-valued convex functions on $\R^d$ are continuous (see e.g. \cite[Corollary 10.1.1]{rockafellar1997convex}). 

% \smallskip
% \noindent 
% \textit{\ref{prop6}} First, notice that $\psi_n\in{\mathcal{C}}_{b,p}(\R^d)$ for all $n\in\mathbb N$. Then by \textit{\ref{prop3}}, $\psi_n\geq \psi$ and $\psi_n\geq \psi_{n+1}$, imply $\conv{\psi_n}\geq \conv{\psi}$ and $\conv{\psi_n}\geq \conv{\psi_{n+1}}$. This ensures that the pointwise limit of $\conv{\psi_n}$ exists and is greater than or equal to $\conv\psi$. So we have
% \begin{align} \label{displ 1}
%     \conv{\psi}(y)\leq \lim_{n\to\infty} \conv{\psi_n}(y)=\lim_{n\to\infty} {}\inf_{\substack{\rho \in\PP_{p}(\R^d)\\ \delta_y\lc\rho}}\rho(\psi_n),
% \end{align}
% where we applied Lemma \ref{repr lemm} to $\psi_n$. But by the same lemma applied to $\psi$, one can find a minimizing sequence $\{\rho_k\}_{k\in\mathbb N} \subseteq \PP_{p}(\R^d)$ with $\delta_y\lc\rho_k$ for all $k\in\mathbb N$ such that $\rho_k(\psi)\to\conv{\psi}(y)$. This implies
% $$ \lim_{n\to\infty} \conv{\psi_n}(y)\leq \lim_{n\to\infty} {}\rho_k(\psi_n)= {}\rho_k(\psi) + \lim_{n\to\infty} {}\frac{\rho_k\left(|\cdot|^p\right)}{n}=\rho_k(\psi), \text{ for all $k\in\mathbb N,$}$$
% where the limit vanishes since each $\rho_k$ has finite $p$-th moment. Letting $k\to\infty$ proves the claim.

% \smallskip

% Properties \textit{\ref{prop1}--\textit{\ref{prop6}}} for the increasing convex hull are proven analogously.
% \end{proof}


% The following is a technical lemma which we shall later require in the proof of the main result. See Appendix \ref{dokaz konv cudo} for the proof.



% The following two lemmas pertain to the work of \text{Ciosmak} \cite{ciosmak2024localisationconstrainedtransportsi,ciosmak2024localisationconstrainedtransportsii}, where the irreducible convex pavings of martingale transport are generalized to irreducible pavings induced by a given convex cone. 

% \begin{lemma}[{\cite[Lemma 4.8]{ciosmak2024localisationconstrainedtransportsi}}] \label{ciomask prva} Let $f\in\bar D$. Then there exists an increasing net $\{f_\alpha\}_{\alpha\in A} \subseteq D$ converging to $f$.
    
% \end{lemma}
% \begin{lemma}[{\cite[Lemma 4.9]{ciosmak2024localisationconstrainedtransportsi}}] \label{ciomak druga} Let $\mu$ be a finite Radon measure on a topological space $\Omega$ and let $\{f_{\alpha}\}_{\alpha \in A}$ be an increasing net of lower semicontinuous functions on $\Omega$ converging to $f$. Assume that for some $\alpha_0\in A$, $f_{\alpha_0}$ and $f$ are $\mu$-integrable. Then
% $$\lim_{\alpha\in A}\mu(f_\alpha)=\mu(f).$$
    
% \end{lemma}
\subsection{Weak Optimal Transport and Duality}\label{WOTD}
We now introduce some key notions related to weak optimal transport. As is common in weak transport, we assume that $C:X\times\PP_p(Y)\to[0,\infty]$ is lsc and the map $\rho\mapsto C(x,\rho)$ is convex for each $x\in X$.
    

% \begin{remark} \label{cplusrho rmk}
%    Note that the cost function related to the stretched Brownian motion (see Subsection \ref{sbm}) fails to satisfy the third point, but the modification $C(x,\rho)+\int1/2|y|^2d\rho(y)$ satisfies all three points. This motivates us to adopt the weaker assumption that $C_{f_0}(x,\rho):=C(x,\rho)+\int f_0d\rho$ satisfies Condition \hyperref[condA]{(A)} for some $f_0\in{\mathcal{C}}_{b,p}(Y)$.
% \end{remark}

% \begin{definition}[$C_{f_0}$]\label{cf0} Given a measurable map  $$C:X\times\PP_p(Y)\to(-\infty,\infty]$$ and  $f_0\in{\mathcal{C}}_{b,p}(Y)$, we define  $$C_{f_0}:X\times\PP_p(Y)\to(-\infty,\infty]$$
% by 
%  $$C_{f_0}(x,\rho):=C(x,\rho)+\int_Y f_0(y)d\rho(y).$$
    
% \end{definition}


% We shall denote by ${\mathcal{C}}_{b,p}(Y)+f_0$ the set of all functions of the form $\phi+f_0$, where $\phi\in{\mathcal{C}}_{b,p}(Y)$; we define ${\mathcal{C}}^{lsc}_{b,p}(Y)+f_0$ analogously.

\begin{definition}[$C$-conjugate] \label{cconj}
   Let $C:X\times\PP_p(Y)\to[0,\infty]$ be measurable. We define the $C$-conjugate $\psi^C:X\to (-\infty,\infty]$ by
    \begin{equation*}
        \psi^C(x):=\inf_{\rho\in\PP_p(Y)} \rho(\psi)+C(x,\rho), \;\;\psi\in{\mathcal{C}}^{}_{b,p}(Y).
    \end{equation*}
\end{definition}
The function $\psi^C$ is lower semianalytic, in particular, universally measurable (see \cite[Proposition 7.47]{bertsekas1978stochastic}). It is also bounded from below, hence $\mu(\psi^C)$ is well-defined for all measures $\mu\in\PP(X)$.


\begin{lemma} \label{rc stabil lemma}
    Let $C:X\times\PP_p(\R^d)\to[0,\infty]$ be measurable and $\psi\in {\mathcal{C}}_{b,p}(\R^d)$. If $C$ is $\lc$-decreasing in the second argument, then we have $$ \psi^C=({\conv{\psi}})^C,$$
    and if $C$ is $\lic$-decreasing in the second argument, then we have $$ \psi^C=({\iconv{\psi}})^C.$$
\end{lemma}

\begin{proof} We prove the first claim, as the second is proven identically.
    Fix $\psi\in{\mathcal{C}}_{b,p}(\R^d)$. Then
    $$ \psi^C(x)=\inf_{\rho\in\PP_p(\R^d)} \rho(\psi)+C(x,\rho)\geq \inf_{\rho\in\PP_p(\R^d)} \rho(\conv\psi)+C(x,\rho)=(\conv{\psi})^C(x),$$
    where the inequality is by Lemma \ref{repr lemm}\textit{\ref{prop1}}. 
    
    The converse inequality is equivalent to the statement
   \begin{equation} \label{nejjed}
        \forall\varepsilon>0\; \forall\rho\in\PP_p(\R^d) \;\exists\widetilde{\rho}\in\PP_p(\R^d):\widetilde{\rho}(\psi)+C(x,\widetilde{\rho})\leq {\rho}(\conv{\psi})+C(x,\rho)+\varepsilon.
   \end{equation}
    Fix $\varepsilon>0$ and $\rho\in\PP_p(\R^d).$ Then by Lemma \ref{repr lemm}\textit{\ref{prop6}} we have
    $\conv{\psi_n}\downarrow \conv{\psi}$, as $n\to\infty$, where $\psi_n:=\psi+\frac{1}{n}|\cdot|^p$. Since $\rho(\conv{\psi_1})<\infty$, we have by monotone convergence 
    $$\rho(\conv{\psi_n})\xrightarrow{n\to\infty}\rho(\conv{\psi}).$$
    So fix $N\in\mathbb N$ sufficiently large such that
   \begin{equation} \label{suff small}
        \rho(\conv{\psi_N})-\rho(\conv \psi)\leq \varepsilon/2.
   \end{equation}
    Now by Lemma \ref{repr lemm}\textit{\ref{conv repr bullet}} applied to $\psi_N$ we have that the set 
    $$A:=\left\{(y,\xi)\in \R^d\times\PP_p(\R^d) \mid \conv{\psi_N}(y)+\varepsilon/2\geq \xi(\psi_N), \;\delta_y \lc\xi\right\}$$
    satisfies $\text{proj}_1A=\R^d$. It is straightforward to verify that $A$ is Borel measurable and by the von Neumann uniformization theorem, we may extract a universally measurable selection $\{\xi_y\}_{y\in \R^d}\subseteq\PP_p(\R^d)$, such that
    \begin{equation} \label{selec prop}
        \conv{\psi_N}(y)+\varepsilon/2\geq \xi_y(\psi_N), \;\delta_y \lc\xi_y.
    \end{equation}
   We define $\widetilde{\rho}$ by $d\widetilde{\rho}(z):=\int_{\R^d} d\xi_y(z)d\rho(y).$ Integrating the inequality in (\ref{selec prop}) with respect to $\rho$ yields
   \begin{equation} \label{bitna}
       \rho(\conv{\psi_N})+\varepsilon/2\geq \widetilde{\rho}(\psi_N).
   \end{equation}
   On the one hand, since $\psi$ is bounded from below we have
$$\widetilde{\rho}(\psi_N)\geq \ell +1/N\widetilde{\rho}(|\cdot|^p),$$
and on the other hand, by Lemma \ref{repr lemm}\textit{\ref{prop4}}, 
$\rho(\conv{\psi_N}) <\infty.$
These two inequalities combined with (\ref{bitna}) yield that $\widetilde{\rho}\in\PP_p(\R^d)$.

Now notice that (\ref{suff small}) and (\ref{bitna}) imply
\begin{equation} \label{bistna}
       \rho(\conv{\psi})+\varepsilon\geq \widetilde{\rho}(\psi_N)\geq\widetilde{\rho}(\psi).
   \end{equation}
This partially shows the inequality in (\ref{nejjed}). To finalize the proof, note that for arbitrary convex $f$, we have
$$\widetilde{\rho}(f)=\int \xi_y(f)d\rho(y)\geq\int \delta_y(f)d\rho(y)=\rho(f),$$
i.e., $\rho\lc\widetilde{\rho}$, which by the monotonicity assumption implies $C(x,\rho)\geq C(x,\widetilde{\rho})$.
\end{proof}

We conclude this section by recalling the general duality theorem for weak optimal transport.
\begin{theorem}[{\cite[{Theorem 1.3}]{veraguas2019existence}}]\label{th:kantorovichduality}
	 Let $\mu\in\PP(X)$, $\nu \in \PP_p(Y)$ and $C:X\times\PP_p(Y)\to[0,\infty]$ be lsc and convex in the second argument. Then
	\begin{align}
		\inf_{\pi\in\Pi(\mu,\nu)} \int C(x,\pi_x)d\mu(x) = \sup_{\psi \in {\mathcal{C}}_{b,p}(Y)} \mu(\psi^C) - \nu(\psi).
	\end{align}
\end{theorem}

% The second modification provides the duality result for a broader class of cost functions.

% \begin{theorem}\label{th:kantorovssssichsdsdsdduality}
% 	Let $C\colon X\times \mathcal P_{p}(Y)\rightarrow \mathbb{R}\cup\{\infty\}$ be a cost function such that $C_{f_0}$ satisfies Condition \hyperref[condA]{(A)} for some $f_0\in{\mathcal{C}}_{b,p}$. Then we have {\color{black}for all} $\mu\in\mathcal P(X)$ and $\nu\in\mathcal P_p(Y)$
% 	\begin{align*}
% 		\inf_{\pi\in\Pi(\mu,\nu)} \int C(x,\pi_x)d\mu(x)= \sup_{\psi \in {\mathcal{C}}^{}_{b,p}+f_0} \mu(R_C\psi) - \nu(\psi)= \sup_{\psi \in {\mathcal{C}}^{lsc}_{b,p}+f_0} \mu(R_C\psi) - \nu(\psi).
% 	\end{align*}
% \end{theorem}
% \begin{proof}
% By applying Theorem \ref{th:kantorovichduality} to ${C_{f_0}}$, one has for all $\mu\in\PP(X)$ and $\nu\in\PP_p(Y)$
%     	\begin{align*}
% 		\inf_{\pi\in\Pi(\mu,\nu)} \int {C_{f_0}}(x,\pi_x)d\mu(x)= \sup_{\psi \in {\mathcal{C}}^{}_{b,p}} \mu(R_{{C_{f_0}}}\psi) - \nu(\psi).
% 	\end{align*}
%     Expanding ${C_{f_0}}$, we get
%     \begin{align*}
% 		\inf_{\pi\in\Pi(\mu,\nu)} \int {C}(x,\pi_x) +\pi_x(f_0)d\mu(x)&= \sup_{\psi \in {\mathcal{C}}^{}_{b,p}} \mu(R_{ C}(\psi +f_0)) - \nu(\psi)\\
%         &=\sup_{\psi \in {\mathcal{C}}^{}_{b,p}} \mu(R_{ C}(\psi +f_0)) - \nu(\psi+f_0)+\nu(f_0)\\
%         &=\sup_{\psi \in {\mathcal{C}}^{}_{b,p}+f_0} \mu(R_{ C}\psi) - \nu(\psi)+\nu(f_0).
% 	\end{align*}
%     But notice that 
%     $$\inf_{\pi\in\Pi(\mu,\nu)} \int {C}(x,\pi_x) +\pi_x(f_0)d\mu(x)=\inf_{\pi\in\Pi(\mu,\nu)} \int {C}(x,\pi_x) d\mu(x)+\nu(f_0).$$
%     Since $\nu(f_0)<\infty$, we subtract it from both sides to get
%     \begin{align*}
% 		\inf_{\pi\in\Pi(\mu,\nu)} \int {C}(x,\pi_x) d\mu(x)=\sup_{\psi \in {\mathcal{C}}^{}_{b,p}+f_0} \mu(R_{ C}\psi) - \nu(\psi).
% 	\end{align*}

%     Analogously, by applying Theorem \ref{th:kantorovssssichduality} to ${C_{f_0}}$, one shows that
%      \begin{align*}
% 		\inf_{\pi\in\Pi(\mu,\nu)} \int {C}(x,\pi_x) d\mu(x)=\sup_{\psi \in {\mathcal{C}}^{lsc}_{b,p}+f_0} \mu(R_{ C}\psi) - \nu(\psi).
% 	\end{align*}
    
% \end{proof}


\section{Proof of the Main Theorems}\label{MT}


We now prove the main theorem on restricting the dual problem to convex functions.






\begin{proof}[Proof of Theorem \ref{mainthm}]
    Assume that $C$ is $\lc$-decreasing. Fix $\mu\in\PP(X)$ and $\nu\in\PP_p(\R^d)$. Then we have
    \begin{align*}
		\inf_{\pi\in\Pi(\mu,\nu)} \int C(x,\pi_x)d\mu(x)&= \sup_{\psi \in {\mathcal{C}}^{}_{b,p}(\R^d)} \mu(\psi^C) - \nu(\psi)\\
        &\leq \sup_{\psi \in {\mathcal{C}}^{}_{b,p}(\R^d)} \mu( (\conv\psi)^C) - \nu(\conv\psi)
	\end{align*}
    where the equality holds by Theorem \ref{th:kantorovichduality} and the inequality holds by Lemma \ref{rc stabil lemma} and Lemma \ref{repr lemm}\textit{\ref{prop1}}. 
But whenever $\psi\in{\mathcal{C}}_{b,p}(\R^d)$, we have $\conv \psi \in {\mathcal{C}}_{b,p}(\R^d)$, so
    \begin{align*}
        &\leq \sup_{\substack{\psi \in {\mathcal{C}}^{}_{b,p}(\R^d) \\ \psi \text{ convex}}} \mu( \psi^C) - \nu(\psi)\\
        &\leq \sup_{\psi \in {\mathcal{C}}^{}_{b,p}(\R^d)} \mu(\psi^C) - \nu(\psi)\\&=\inf_{\pi\in\Pi(\mu,\nu)} \int C(x,\pi_x)d\mu(x).
	\end{align*}

    For the converse, fix $x\in X$ and $\rho_1,\rho_2\in\PP_p(\R^d)$, such that $\rho_1\lc\rho_2$. By assumption, we have
    \begin{align*} 
       C(x,\rho_1)&=\inf_{\pi\in\Pi(\delta_x,\rho_1)}\int C(\widetilde x,\pi_{\widetilde{x}}) d\delta_x(\widetilde{x})= \sup_{\substack{\psi \in  {\mathcal{C}}^{}_{b,p}(\R^d) \\\psi \text{ convex}}} \psi^C(x) - \rho_1(\psi)\\
       &\geq \sup_{\substack{\psi \in  {\mathcal{C}}^{}_{b,p}(\R^d) \\ \psi\text{ convex}}} \psi^C(x) - \rho_2(\psi)=C(x,\rho_2).
    \end{align*}
\end{proof}
% \fbox{\begin{minipage}{\textwidth}
% \color{black}
% \textbf{Question:} I believe that one should be able to get
% \begin{align*} 
% 		\inf_{\pi\in\Pi(\mu,\nu)} \int C(x,\pi_x)d\mu(x)= \sup_{\color{red}{\psi \in {D}\cap {\mathcal{C}}^{}_{b,p}}} \mu(\psi^C) - \nu(\psi),
% 	\end{align*}
% instead of (\ref{mainthmdispl}). This would be a much more elegant result. This is equivalent to showing,
% \begin{align*} 
% 		 \sup_{\psi \in \bar {D}\cap {\mathcal{C}}^{lsc}_{b,p}} \mu(\psi^C) - \nu(\psi)\leq \sup_{\psi \in {D}\cap {\mathcal{C}}^{}_{b,p}} \mu(\psi^C) - \nu(\psi),
% 	\end{align*}
%     or equivalently showing:
%     $$\forall\varepsilon>0\;\forall \psi\in\bar D\cap {\mathcal{C}}^{lsc}_{b,p}\;\exists\widetilde \psi\in D\cap{\mathcal{C}}_{b,p}: \mu(\widetilde\psi^C)-\nu(\widetilde \psi)\geq\mu(\psi^C)-\nu(\psi)-\varepsilon.$$
% By Lemma \ref{konv cudo}, one can easily construct $\widetilde \psi$ such that $$-\nu(\widetilde \psi)\geq-\nu(\psi)-\varepsilon/2.$$ This approach fails to show the $C$-conjugate part of the inequality, as it would require some regularity. But perhaps a different approximation argument works?
% \end{minipage}}
The increasing convex version of the restriction theorem also holds and is proven analogously.
\begin{theorem} \label{mainthmic}
    Let $\mu\in\PP(X)$, $\nu \in \PP_p(\R^d)$ and $C:X\times\PP_p(\R^d)\to[0,\infty]$ be lsc and convex in the second argument. If $C$ is $\lic$-decreasing in the second argument, then 
\begin{align} \label{mainthmdisplic}
		\inf_{\pi\in\Pi(\mu,\nu)} \int C(x,\pi_x)d\mu(x)= \sup_{\substack{\psi \in  {\mathcal{C}}^{}_{b,p}(\R^d) \\ \psi\textnormal{ inc. convex}}} \mu(\psi^C) - \nu(\psi).
	\end{align}
    Conversely, if (\ref{mainthmdisplic}) holds for all $\mu\in\PP(X)$, $\nu\in\PP_p(\R^d)$, then $C$ is $\lic$-decreasing in the second argument.
\end{theorem}



    The proof of Theorem \ref{mainthm} (and the auxiliary results it uses) can be modified to accommodate cost functions $C$ that take negative values, for example, the martingale Benamou--Brenier problem \cite{backhoffveraguas2025existencebassmartingalesmartingale}. This relaxation is formulated in the following theorem.
    

\begin{theorem} \label{thmmod1}
   Let $\mu\in\PP(X)$, $\nu \in \PP_p(\R^d)$, and $C:X\times\PP_p(\R^d)\to(-\infty,\infty]$ be lsc, convex in the second argument, and such that it satisfies the lower bound $C(x,\rho)\geq-(a_\ell(x)+\rho(b_\ell))$, for some $a_\ell\in L^1(\mu)\cap \LSC(X)$\footnote{$\LSC(X)$ denotes the real-valued lsc functions on $X$} and some convex function $b_\ell\in{\mathcal{C}}_{b,p}(\R^d)$. If $C$ is $\lc$-decreasing in the second argument, then 
    \begin{align*} 
		\inf_{\pi\in\Pi(\mu,\nu)} \int C(x,\pi_x)d\mu(x)= \sup_{\substack{{\psi \in ({\mathcal{C}}_{b,p}(\R^d)+b_\ell)}\\ \psi \textnormal{ convex}}} \mu(\psi^C) - \nu(\psi).
	\end{align*}
\end{theorem}
The analogous result also holds for the increasing convex case.

\section{Attainment considerations}
\label{sec att}
While duality holds in general, the supremum is only attained under additional regularity assumptions.  Dealing with the delicate question of dual attainment requires us to define the dual problem appropriately, so we follow the treatment in \cite{BePaRiSc25}. Let $X,Y$ be Polish spaces, $\mu\in\P(X)$ and $\nu\in\PP_p(Y)$. Let us denote by $\mathcal L^1(\rho)$ the space of $(-\infty,\infty]$-valued Borel functions that are $\rho$-integrable. 

 

To guarantee attainment, we shall require two conditions. Firstly, in line with the classical theory, we require a \textit{boundedness condition}
\begin{equation*}
    \label{condB} \tag{B}
    C(x,\rho)\leq a(x)+\rho(b)
\end{equation*}
for a $\mu$-integrable function $a:X\to\R$ and a $\nu$-integrable function $b:Y\to\R$.

Secondly, we require a mild \textit{continuity condition}, that is, for any increasing sequence $(Y_k)_k$ of Borel sets with $\bigcup_kY_k=Y$, we have 
\begin{equation*}
    \label{condC} \tag{C}
    C(x,\rho)\geq \limsup_k C(x,(\rho|Y_k)/\rho(Y_k)).
\end{equation*}

Let us emphasize the subtle difference between Definition \ref{cconj} and the way we define the $C$-conjugate in this section. For a measurable function $\psi:Y\to(-\infty,\infty]$, we have
$$\psi^C(x):=\inf_{\substack{\rho\in\PP_p(Y)\\ \text{s.t. }\psi\in\mathcal L ^1(\rho)} } \rho(\psi)+C(x,\rho),$$
where the infimum over an empty set is interpreted as $+\infty$. In the case of $\psi\in{\mathcal{C}}_{b,p}(Y)$, the two definitions coincide, but for less regular functions they may differ.

Lastly, to precisely state the dual formulation, we define the set of admissible pairs $\Phi_C(\mu,\nu)$ as pairs of functions $(\phi,\psi)\in (-\mathcal L^1(\mu))\times \mathcal L^1 (\nu)$ satisfying 
$$\phi(x)-\rho(\psi)\leq C(x,\rho), $$
for all $x\in X$ and $\rho\in\P_p(Y)$ such that $\psi\in \mathcal L^1(\rho)$.

We can now state the main attainment result. 

\begin{theorem} \label{thmatt}
    Let $X$ be a Polish space, $\mu\in\PP(X)$ and $\nu\in\PP_p(\R^d)$. Assume that $C:X\times\PP_p(\R^d)\to[0,\infty]$ is lsc and both convex and  $\lc$-decreasing in the second argument.  If $C$ satisfies conditions (\ref{condB}) and (\ref{condC}), then
   \begin{equation}
        \inf_{\pi\in\Pi(\mu,\nu)}\int C(x,\pi_x)d\mu(x)=\sup_{\substack{(\psi^C,\psi)\in\Phi_C(\mu,\nu) \\ \psi \textnormal{ convex}}} \mu(\psi^C)-\nu(\psi) \label{att_eq}
   \end{equation}
and the supremum is attained. 
\end{theorem}
\begin{proof}
    The equality in (\ref{att_eq}) is an immediate consequence of \cite[Theorem 2.2]{BePaRiSc25} and Theorem \ref{mainthm}.

We now proceed to show that the supremum is attained. Let us call $(\phi,\psi)\in\Phi_C(\mu,\nu)$ an \textit{optimal dual pair} if $$\mu(\phi)-\nu(\psi)=\sup_{\substack{(\phi,\psi)\in\Phi_C(\mu,\nu)}} \mu(\phi)-\nu(\psi).$$
By \cite[Theorem 2.2]{BePaRiSc25}, an optimal dual pair $(\phi,\psi)\in\Phi_C(\mu,\nu)$ exists.
Assume that  $((\conv\psi)^C,\conv \psi)$ is also an optimal dual pair. Then:
$$\mu((\conv\psi)^C)-\nu(\conv \psi)=\sup_{\substack{(\phi,\psi)\in\Phi_C(\mu,\nu)}} \mu(\phi)-\nu(\psi)\geq\sup_{\substack{(\psi^C,\psi)\in\Phi_C(\mu,\nu)\\\psi \textnormal{ convex}}} \mu(\psi^C)-\nu(\psi),$$
which implies that the supremum in (\ref{att_eq}) is attained.


We are left to verify that $((\conv\psi)^C,\conv \psi)$ is an optimal dual pair. We do this in three steps. First, we show that $(\phi,\convr\psi)$ is an optimal dual pair, where $\convr$ is an appropriately defined approximation of the convex hull operator. Then we pass to the limit to show that $(\phi,\conv\psi)$ is also an optimal dual pair. Finally, we show that $\phi$ may be replaced by $(\conv\psi)^C$ with no ill effect on the optimality. 

\textit{Step 1:} $(\phi,\convr\psi) $ \textit{is an optimal dual pair}

\noindent Let us first define the operator $\convr$. For $R\in \mathbb N$, let $$\convr\psi(y):=\inf\left\{ \xi(\psi)\midauto \xi\in\P(B_R(y)),\, \delta_y\lc\xi,\, |\textnormal {supp}(\xi)|\leq d+1\right\}.$$ 
Compare this with the following representation of the convex hull:
\begin{equation}
\label{conv_hull_gen}
\conv\psi(y)=\inf\left\{ \xi(\psi)\midauto \xi\in\P(\R^d),\, \delta_y\lc\xi,\, |\textnormal {supp}(\xi)|\leq d+1\right\}.
\end{equation} 
%This representation of $\conv\psi$ is valid for arbitrary $(-\infty,\infty]$-valued functions $\psi$ on $\R^d$ (see \cite[Corollary 17.1.5]{rockafellar1997convex}).

Note that $\convr\psi$ is \textit{not} necessarily a convex function, but one has $\conv\psi \leq\convr\psi\leq\psi$ and $\convr\psi\downarrow\conv\psi$ as $R\to\infty$. Since $\convr\psi\leq\psi$, optimality of $(\phi,\convr\psi)$ will follow as soon as we show that $(\phi,\convr\psi) \in\Phi_C(\mu,\nu).$

Let us show that $\convr\psi \in \mathcal L^1(\nu)$. Assume that $\convr\psi(y)=-\infty$ for some $y\in \R^d$. This implies that $\conv\psi(y)=-\infty$ and thus by (\ref{conv_hull_gen}) we can find a sequence of probability measures $\{\xi_k\}_{k=1}^\infty$, each supported on at most $d+1$ points with mean $y$, such that $\lim_k \xi_k(\psi)=-\infty$. Fix $x_0\in X$ such that $\phi(x_0)\in\R$. Then we use the fact that $(\phi,\psi)\in\Phi_C(\mu,\nu)$ and the $\lc$-monotonicity of $C$ to get
$$\phi(x_0)-\xi_k(\psi)\leq C(x_0,\xi_k)\leq C(x_0,\delta_y)\leq a(x_0)+b(y).$$
Then the left-hand side tends to $\infty$, but the right-hand side is finite, leading to a contradiction. Thus we have shown that $\convr\psi$ is $(-\infty,\infty]$-valued. Note that our proof also shows that $\conv\psi$ is $(-\infty,\infty]$-valued, and hence a proper convex function. By \cite[Corollary 12.1.2]{rockafellar1997convex}, $\conv \psi$ admits an affine minorant  $\textnormal{aff}:\R^d\to\mathbb R$, so
$$\textnormal{aff}(y)\leq\conv\psi(y)\leq \convr\psi(y)\leq\psi(y),$$
but since $\psi$ and $\textnormal{aff}$ are both $\nu$-integrable, so is $\convr\psi$ (and $\conv \psi$).

Let us now show that the defining inequality of $\Phi_C(\mu,\nu)$ is satisfied by $(\phi,\convr\psi)$. Fix $\varepsilon>0$, $x\in X$ such that $\phi(x)\in\R$ and $\rho \in \P_p(\R^d)$ with $\convr\psi \in \mathcal L^1(\rho)$. By the von Neumann uniformization theorem, we can extract a measurable collection of probability measures $\{\xi_y\}_{y\in\R^d}\subseteq\P_p(\R^d)$ such that \begin{equation}
    \label{meas_sel_ineq}\convr\psi(y)\geq\xi_y(\psi)-\varepsilon, \;\, \delta_y\lc\xi_y, \;\,\xi_y(B_R(y))=1.
\end{equation}
Define $d\widetilde{\rho}(z):=d\xi_y(z) d\rho(y).$ Then $\widetilde{\rho}\ \in \P_p(\R^d)$ since $\widetilde{\rho}(|\cdot|^p)\leq2^{p-1}[{\rho}(|\cdot|^p)+R^p]$. Integrate the inequality in (\ref{meas_sel_ineq}) with respect to $\rho$ to get 
\begin{equation}
\label{ineq_meas_se}
\rho(\convr\psi)\geq\widetilde\rho(\psi)- \varepsilon.
\end{equation}
Note that $\psi\in\mathcal L^1(\widetilde\rho)$ since $\widetilde{\rho}(|\psi|)\leq \rho(\convr\psi)+\varepsilon+2\widetilde{\rho}(\textnormal{aff}^-)<\infty.$ 

Thus, we have
$$\phi(x)-\rho(\convr\psi)\leq \phi(x)-\widetilde\rho(\psi)+\varepsilon\leq C(x,\widetilde{\rho})+\varepsilon\leq C(x,\rho)+\varepsilon,$$
where the inequalities follow from (\ref{ineq_meas_se}), the fact that $(\phi,\psi)\in \Phi_C(\mu,\nu)$ and $\lc$-monotonicity of $C$, respectively. 

This shows that $(\phi,\convr\psi)\in\Phi_C(\mu,\nu)$, which immediately implies that $(\phi,\convr\psi)$ is an optimal dual pair.


\textit{Step 2:} $(\phi,\conv\psi) $ \textit{is an optimal dual pair}

\noindent
As before, since $\conv\psi\leq\psi$, optimality will follow as soon as $(\phi,\conv\psi)\in\Phi_C(\mu,\nu)$. We already verified that $\conv\psi\in\mathcal L^1(\nu)$ in the previous step, so it remains to show that the defining inequality of $\Phi_C(\mu,\nu)$ holds for the pair $(\phi,\conv\psi)$.

Fix $x\in X$ such that $\phi(x)\in\R$ and $\rho \in \P_p(\R^d)$ with $\conv\psi \in \mathcal L^1(\rho)$. By Egorov's theorem, we can find a sequence $\{A_k\}^{\infty}_{k=1}$ of increasing compact subsets of $\R^d$ such that $\rho(\bigcup_kA_k)=1$ and $\convr\psi\to\conv\psi$ uniformly on each set $A_k$. Define for each $k\in\mathbb N$ a probability measure $\rho_k:=(\rho|{A_k}) / \rho(A_k).$

Then by dominated convergence, we have
$$\phi(x)-\rho(\conv\psi)=\lim_k\phi(x)-\rho_k(\conv\psi),$$
but by uniform convergence we also get
$$\lim_k\phi(x)-\rho_k(\conv\psi)=\lim_k \lim_R\phi(x)-\rho_k(\convr\psi)\leq \limsup_k C(x,\rho_k)\leq C(x,\rho),$$
where the inequalities above are due to $(\phi,\convr\psi)\in\Phi_C(\mu,\nu)$ and the continuity condition (\ref{condC}) with $Y_k:=A_k\cup(\R^d\setminus\bigcup_{k=1}^\infty A_k)$.

\textit{Step 3:} $((\conv\psi)^C,\conv\psi) $ \textit{is an optimal dual pair}

\noindent
Since $(\phi,\conv\psi)\in\Phi_C(\mu,\nu)$, we have for all $x\in X$ and $\rho \in \P_p(\R^d)$ such that $\conv\psi \in\mathcal L^1(\rho)$ 
$$\phi(x)\leq \rho(\conv\psi)+C(x,\rho),$$
and taking the infimum over all such $\rho$ yields $\phi(x)\leq (\conv \psi)^C(x)$. On the other hand, by the definition of the convex conjugate, we have 
$$(\conv\psi)^C(x)\leq \rho(\conv\psi)+C(x,\rho)$$
for all $x\in X$ and $\rho\in\P_p(\R^d)$ such that $\conv \psi \in \mathcal L^1(\rho)$. %By taking $\rho=\delta_{y_0}$ for $ y_0\in\R^d$ such that $\conv\psi(y_0)\in\R$, we see that $(\conv\psi)^C$ is $[-\infty,\infty)$-valued. 

By $\nu$-integrability of $\conv\psi$  and assumption (\ref{condB}), we get
$$\phi(x)\leq(\conv\psi)^C(x)\leq \nu(\conv\psi)+a(x)+\nu(b),$$
 which proves that $(\conv\psi)^C\in(-\mathcal L^1(\mu))$. 
 
 Hence $((\conv\psi)^C,\conv\psi)\in \Phi_C(\mu,\nu)$ and since $\phi\leq (\conv\psi)^C$, we have that $((\conv\psi)^C,\conv \psi)$ is an optimal dual pair.


 
\end{proof}
The proof given above, \textit{mutatis mutandis}, also shows the increasing convex version of Theorem \ref{thmatt}.

\begin{theorem} \label{thmatt2}
    Let $X$ be a Polish space, $\mu\in\PP(X)$ and $\nu\in\PP_p(\R^d)$. Assume that $C:X\times\PP_p(\R^d)\to[0,\infty]$ is lsc and both convex and  $\lic$-decreasing in the second argument. If $C$ satisfies conditions (\ref{condB}) and (\ref{condC}), then
   \begin{equation}
        \inf_{\pi\in\Pi(\mu,\nu)}\int C(x,\pi_x)d\mu(x)=\sup_{\substack{(\psi^C,\psi)\in\Phi_C(\mu,\nu) \\ \psi \textnormal{ increasing convex}}} \mu(\psi^C)-\nu(\psi)\label{att_eq2}
   \end{equation}
   and the supremum is attained. 
\end{theorem}

% \begin{remark}
%     In contrast to the dual formulation of Theorem \ref{mainthm}, we generally do not expect dual potentials to be continuous, but rather only $\nu$-integrable. However, in certain less regular settings a nonintegrable convex dual potential could exist, see Remark \ref{nonintdual}.
% \end{remark}

\begin{remark} \label{lowerboundremark}
    The assumption that $C$ is nonnegative in Theorem \ref{thmatt} can be weakened. Indeed, as in Theorem \ref{thmmod1}, it suffices to have $C(x,\rho)\geq-(a_{\ell}(x)+\rho(b_{\ell}))$ for $a_{\ell}\in L^1(\mu)\cap \LSC(X)$ and some convex function $b_{\ell}\in\mathcal{C}_{b,p}(\R^d)$. 
\end{remark}



\section{Applications}\label{APP}

\subsection{Convex Dual Potentials}




\subsubsection{Barycentric Costs}


\label{barcos}
A cost function that depends on the second argument solely via its mean/\textit{barycenter}, is called barycentric.
In \cite{GoJu18}, \text{Gozlan} and \text{Juillet} consider the dual formulation associated to such costs and restrict it to convex functions. We recover this restriction result using Theorem \ref{mainthm}.

Let $X=\R^d$ and $p=1$, and take $\mu,\nu\in\P_1(\R^d).$ We consider the cost function $C(x,\rho)=\theta(x-\mean(\rho))$ for some convex $\theta: \R ^d \to [0,\infty)$, where $\mean(\rho):=\int_{\R^d} y d\rho(y)$. Note that $\rho_1\lc\rho_2$ implies $\mean(\rho_1)=\mean(\rho_2)$, which shows that $C$ is $\lc$-decreasing, so by Theorem \ref{mainthm}, we get
$$\inf_{\pi\in\Pi(\mu,\nu)} \int \theta(x-\mean(\pi_x))d\mu(x)= \sup_{\substack{\psi \in {\mathcal{C}}^{}_{b,1}(\R^d) \\ \psi\text{ convex}}} \mu(\psi^C) - \nu(\psi).$$
This recovers {\cite[Theorem 1.1]{GoJu18}}. It is worth noting that the $C$-conjugate in the expression above simplifies to $\psi^C(x)=
\inf_{z\in\R^d}\psi(z)+\theta(x-z).$

For concreteness, taking $\theta(x)=\left \| x \right \|$, for $\left \| \cdot \right \|$ a norm on $\R^d$, produces the convex Kantorovich--Rubinstein duality, while $\theta(x)=|x|^2$ with $\mu,\nu\in\PP_2(\R^d)$ recovers the Brenier--Strassen mixture result \cite[Theorem 1.2]{GoJu18}. Moreover, Theorem \ref{thmatt} ensures that in both of these instances the dual is attained by a $\nu$-integrable convex function $\psi^{\text{opt}}:\R^d\to(-\infty,\infty]$.

\subsubsection{Martingale Benamou--Brenier} \label{sbm}
The martingale analogue of the Benamou--Brenier problem \cite{backhoffveraguas2025existencebassmartingalesmartingale}  turns out to be equivalent to a weak transport problem where the dual formulation can be restricted to convex functions. We reproduce this result using Theorem \ref{thmmod1}.

Let $X=\R^d$, $p=2$, $a_\ell(x)=0$ and $b_\ell(y)=1/2|y|^2\in{\mathcal{C}}_{b,2}(\R^d)$. We consider the following cost function:
$$C(x,\rho)=  \begin{cases} 
        -\text{MCov}(\rho,\gamma), & x=\mean (\rho) \\
      \infty, & \text{otherwise}, 
   \end{cases}
  $$
where 
$\text{MCov}(\rho_1,\rho_2):=\max_{q \in \Pi(\rho_1,\rho_2) } \int_{\R^d\times \R ^d}yz dq(y,z)$ and $\gamma$ denotes the standard normal distribution on $\R^d$.
 
Let $\mu,\nu\in\PP_2(\R^d)$ with $\mu\lc\nu$. One readily checks the assumptions of Theorem \ref{thmmod1}, hence we have
$$\sup_{\pi\in\text{MT}(\mu,\nu)} \int \text{MCov}(\pi_x,\gamma)d\mu(x)= \inf_{\substack{{\psi \in \mathcal Q} \\\psi \text{ convex} }} \nu(\psi) -\mu(\psi^C),$$
where $\text{MT}(\mu,\nu)$ denotes the set of martingale couplings between $\mu$ and $\nu$, and $\mathcal{Q}:=( {\mathcal{C}}^{}_{b,2}(\R^d)+1/2|\cdot|^2)$. Thus we have recovered \cite[Proposition 3.5]{backhoffveraguas2025existencebassmartingalesmartingale}.

\begin{remark} \label{nonintdual}
    We note that dual attainment is a subtle issue in the case of martingale transport, which in general does not satisfy the boundedness assumption (\ref{condB}). See \cite{backhoffveraguas2025existencebassmartingalesmartingale} for a detailed study of dual attainment for the martingale Benamou--Brenier problem and, in particular, \cite[Example 6.7]{backhoffveraguas2025existencebassmartingalesmartingale} for an example where there is no optimal $\nu$-integrable dual potential.
    \end{remark}
    \begin{remark} In \cite[Theorem 5.4]{BePaRiSc25}, the authors relax the martingale constraint by considering the weak transport problem with the cost function $C(x,\rho):=\beta|x-\mean (\rho)|^2-\alpha\text{MCov}(\rho,\gamma)$ for $\alpha,\beta>0.$ In this relaxed setting, dual attainment follows from Theorem \ref{thmatt} and Remark \ref{lowerboundremark}. Note that letting $\beta\to\infty$ formally recovers the martingale constraint.
\end{remark}


\subsubsection{Classical Optimal Transport} We highlight the connection with classical transport and apply our results to the celebrated quadratic cost function. 

Let $X=\R^d$,  $p=2$, $a_\ell(x)=\frac{x^2}{2} , b_\ell(y)=\frac{y^2}{2}$ and take $\mu,\nu\in\P_2(\R^d)$. We consider the cost function $C(x,\rho)=\int-xyd\rho(y)$. By Theorem \ref{thmmod1} we have
$$\sup_{\pi\in\Pi(\mu,\nu)}\int xy d\pi(x,y)=\inf_{\substack{\psi\in\mathcal Q\\\psi \text{ convex}}}\mu(\psi^*)+\nu(\psi),$$
where $\psi^*(x):=\sup_{y\in\R^d}xy-\psi(y)$ denotes the convex conjugate of $\psi$. This recovers the classical duality in \cite[Theorem 5.10]{villani2008optimal} for $c(x,y)=-xy.$ The dual is attained by a $\nu$-integrable convex function $\psi^{\text{opt}}:\R^d\to(-\infty,\infty]$. 

\begin{remark}
    More generally, in the context of classical transport---when $C(x,\rho)=\int c(x,y)d\rho(y)$---the condition that $C$ is $\lc$-decreasing in the second argument is equivalent to $c$ being concave in the second argument.
\end{remark}
\subsubsection{Strassen's Theorem for Martingales}
Recall the classical result \cite{Strassen1965} of Strassen on the existence of martingale couplings.
\begin{theorem}[Strassen] \label{strass} Let $\mu,\nu\in\PP_1(\R^d)$. Then $\mu\lc\nu$ if and only if there exists a martingale coupling between $\mu$ and $\nu$, that is, for some $\pi\in\Pi(\mu,\nu)$ we have $x=\textnormal{mean}(\pi_x)$ for $\mu$-a.e. $x\in \R^d$.
    
\end{theorem}
\begin{proof}
    By applying Theorem \ref{mainthm} to $$C(x,\rho):=\begin{cases}
        0 \quad\;\; \delta_x \lc\rho,\\
        \infty \quad \text{otherwise,}
    \end{cases}$$
we get

$$\inf_{\pi\in\Pi(\mu,\nu)} \int \chi_{\delta_x\lc\pi_x} d\mu(x)= \sup_{\substack{\psi \in {\mathcal{C}}^{}_{b,1}(\R^d) \\ \psi\text{ convex}}} \mu(\psi) - \nu(\psi).$$
    The left-hand side is zero exactly when a martingale coupling exists, whereas the right-hand side is zero exactly when $\mu\lc\nu$. This proves the claim.
\end{proof}

Although this strategy for proving Strassen's theorem is identical to that of \cite[Theorem 3.1]{GoRoSaTe14}, the novel contribution lies in the fact that it suffices to check the assumptions of Theorem \ref{mainthm} to ensure that the dual can be taken over convex functions. 

\subsection{Increasing Convex Dual Potentials}


\subsubsection{Multiple-Good Monopolist} \label{monopol}
Originally discussed in \cite{DaDeTz17} by \text{Daskalakis, Deckelbaum,} and\text{ Tzamos}, and later in
\cite{backhoffveraguas2020applicationsweaktransporttheory} by \text{Backhoff} and \text{Pammer}, the multiple-good monopolist problem was shown to admit a weak optimal transport formulation. The result \cite[Theorem 6.1]{backhoffveraguas2020applicationsweaktransporttheory} shows that the dual problem can be restricted to convex and coordinatewise increasing functions. We recover this result as a special case of Theorem \ref{mainthmic}.

Let $X=\R^d$ and fix $\theta\in{\mathcal{C}}_{b,{\color{black}p}}(\R^d)$ convex and take $\mu,\nu\in\P_p(\R^d)$.  We consider  $$C(x,\rho):=\inf_{\xi\lic\rho}\theta(x-\mean(\xi)).$$ 
The function $C$ is lsc, convex in the second argument, bounded from below and, crucially,  $\lic$-decreasing. By Theorem \ref{mainthmic}  we have

$$\inf_{\pi\in\Pi(\mu,\nu)} \int \inf_{\xi\lic\pi_x}\theta(x-\mean(\xi))d\mu(x)= \sup_{\substack{\psi \in {\mathcal{C}}^{}_{b,p}(\R^d) \\ \psi\text{ inc. convex}}} \mu(\psi^C) - \nu(\psi),$$
which recovers \cite[Theorem 6.1]{backhoffveraguas2020applicationsweaktransporttheory}. 

Moreover, by Theorem \ref{thmatt2}, we have that the dual is attained by a $\nu$-integrable increasing convex function $\psi^{\text{opt}}:\R^d\to(-\infty,\infty]$. %This attainment result appears to be novel.

\subsubsection{Increasing Convex Kantorovich--Rubinstein}
Let $X=\R^d$,  $\|\cdot\|_q$ be the $\ell^q$ norm on $\R^d$ for $q\in[1,\infty]$. Consider the cost function $C(x,\rho):=\|(x-\mean(\rho))_+\|_q$, where $(v)_+$ denotes the coordinatewise positive part of $v\in \R^d$. By Theorem \ref{mainthmic}, we have for all $\mu,\nu\in\PP_1(\R^d)$:
$$\inf_{\pi\in\Pi(\mu,\nu)} \int \|(x-\mean(\pi_x))_+\|_q d\mu(x)=\sup_{\substack{\psi \; \|\cdot\|_q\text{-Lipschitz}\\\psi\text{ inc. convex}}} \mu(\psi^C)-\nu(\psi).$$
One then replaces $\psi$ by $\psi^C$, which is $1$-Lipschitz and remains increasing and convex. So up to renaming we get
$$\inf_{\pi\in\Pi(\mu,\nu)} \int \|(x-\mean(\pi_x))_+\|_q d\mu(x)=\sup_{\substack{\psi \;1\text{-}\| \cdot\|_q\text{-Lipschitz}\\\psi\text{ inc. convex}}} \mu(\psi)-\nu(\psi).$$
By Theorem \ref{thmatt2}, the dual problem is attained by an increasing convex function $\psi^{\textnormal {opt}}:\R^d\to(-\infty, \infty ]$.


\subsubsection{Strassen's Theorem for Submartingales}
In the same spirit as Theorem \ref{strass}, we recover Strassen's theorem for submartingales as a direct consequence of our main result for increasing convex functions.
\begin{theorem}[Strassen] \label{strass2} Let $\mu,\nu\in\PP_1(\R^d)$. Then $\mu\lic\nu$ if and only if there exists a submartingale coupling between $\mu$ and $\nu$, that is, for some $\pi\in\Pi(\mu,\nu)$ we have $x\leq\textnormal{mean}(\pi_x)$ for $\mu$-a.e. $x\in \R^d$.
    
\end{theorem}





 \section{General Result}\label{nesto}


The final section is devoted to developing a general framework for the restriction of the dual problem. Let $(X,d_X)$ and $(Y,d_Y)$ be Polish metric spaces and $p\geq1$. We write $\LSC_p(Y)$ for the space of lsc functions on $Y$ that are bounded in absolute value by a multiple of $1+d_Y(\cdot,y_0)^p$ for some $y_0\in Y$.

   We call $E\subseteq \LSC_{p}(Y)$ a \emph{stable cone of functions}  if
\begin{enumerate}
    \item \label{ass1} $E$ contains all constant functions,
    \item \label{ass2} $E$ is a convex cone,
    \item \label{ass3} if $\{f_i\}_I\subseteq E$ is such that $\sup_I f_i \in \LSC_p(Y)$, then $\sup_I f_i\in E$.
    \item \label{ass4} for any $f\in E$, there is a collection of \textit{continuous} functions $\{f_i\}_I\subseteq E$ such that $f=\sup_I f_i$,
    \item  \label{ass5} we have $\sup\{\phi(y)\mid \phi\leq f, \phi\in E\}=\inf\{\rho(f)\mid\delta_y\ld \rho \}$ for $f\in{\mathcal{C}}_{b,p}(Y)$.
\end{enumerate}
We say that $\rho_1, \rho_2 \in\PP_p(Y)$ are in $E$-order, in signs $\rho_1\preceq_E \rho_2$, if $\rho_1(f) \leq \rho_2(f)$ for all $f\in E$. 

The proof of Theorem \ref{mainthm} naturally extends to the setting of stable cones.
\begin{theorem}\label{mainthmgen}
     Let $E$ be a stable cone of functions on $Y$, $\mu\in\PP(X)$, $\nu \in \PP_p(Y)$, and $C:X\times\PP_p(Y)\to[0,\infty]$ be lsc and convex in the second argument. If $C$ is $\preceq_E$-decreasing in the second argument, then 
\begin{align} \label{mainthmdisplgen}
		\inf_{\pi\in\Pi(\mu,\nu)} \int C(x,\pi_x)\, d\mu(x)= \sup_{\substack{\psi \in E\\ \psi \textnormal{ lower bdd.}}} \mu(\psi^C) - \nu(\psi).
	\end{align}
    Conversely, if (\ref{mainthmdisplgen}) holds for all $\mu\in\PP(X)$, $\nu\in\PP_p(Y)$, then $C$ is $\preceq_E$-decreasing in the second argument.
\end{theorem}


% \begin{remark}
%     The set $E_b$ denotes all the functions in $E$ that are bounded from below.
% \end{remark}
\begin{remark}
    Note that by taking $E=\{f\in\LSC_p(\R^d): f\text{ convex}\}$ and $E=\{f\in\LSC_p(\R^d): f\text{ increasing convex}\}$ we recover Theorem \ref{mainthm} and Theorem \ref{mainthmic}.
\end{remark} 

We say that $E\subseteq \LSC_p(Y)$ is a \textit{semi-stable cone} if it satisfies assumptions (\ref{ass1})--(\ref{ass4}); these are natural requirements reflecting the basic properties of (increasing) convex functions. The somewhat artificial assumption (\ref{ass5}) is a statement about the dual representation of the $E$-hull, a generalization of the convex hull. It is classical that both the convex hull and the increasing convex hull satisfy such a representation.  Furthermore, if $Y$ is a \textit{compact} Polish space, and $E$ is a semi-stable cone, then by \cite[Corollary XI, T46] {meyer1966probability} it is a stable cone. Despite \cite{Gogus2013} asserting that the same conclusion remains true when  $Y$ is \textit{locally compact}, \cite{Khabibullin2021} points out a flaw in the argument of \cite{Gogus2013}, and \cite{Djire_Wiegerinck_2019} provides a counterexample. 

% Thus, finding appropriate conditions to ensure assumption (\ref{ass5}) holds in the non-compact setting remains an open problem. 


% As was mentioned in the introduction, using essentially the same steps utilized in , we establish the equivalence principle for a general Polish space with respect to a \text{stable cone of functions}.





% \begin{example} The collection of bounded lsc (increasing) convex functions on any closed ball in $\R^d$ forms a stable cone of functions. 
% \end{example}

% \begin{theorem}\label{mainthmintro123}
%     Let $X,Y$ be Polish spaces with $Y$ compact, $\mu\in\PP(X)$, $\nu \in \PP(Y)$  and let $C:X\times\PP(Y)\to[0,\infty]$ be lsc and convex in the second argument. Let $E$ be a stable cone of functions. If $C(x,.)$ is $\preceq_E$-decreasing for every $x\in X$, then 
% \begin{align} \label{mainthmdisplintro123}
% 		\inf_{\pi\in\Pi(\mu,\nu)} \int C(x,\pi_x)\, d\mu(x)= \sup_{\psi \in E} \mu(\psi^C) - \nu(\psi).
% 	\end{align}
%     Conversely, if (\ref{mainthmdisplintro123}) holds for all $\mu\in\PP(X)$, $\nu\in\PP(Y)$, then $C$ is $\preceq_E$-decreasing.
% \end{theorem}




% If $Y$ is a compact space, and $E$ is only assumed to satisfy assumptions (\ref{ass1})--(\ref{ass4}) of the definition of a stable cone, then assumption (\ref{ass5}) is automatically satisfied as well (see \cite[Corollary XI, T46] {meyer1966probability}).

In the compact setting, we  recover the generalization of Strassen's theorem to convex cones; these particular couplings are sometimes referred to as \textit{dilations} (see \cite[Theorem 19.40]{aliprantis2007infinite}).

\begin{theorem}[Dilations]\label{dil} Assume that $Y$ is compact and  $E$ is a semi-stable cone of functions on $Y$. Let $\mu,\nu\in\PP(Y)$. Then $\mu\ld\nu$ if and only if there exists a coupling $\pi\in\Pi(\mu,\nu)$ satisfying $\delta_y \ld\pi_y$ for $\mu$-a.e. $y\in Y$.
\end{theorem}

\begin{proof}
    Apply Theorem \ref{mainthmgen} to $C:Y\times\P(Y)\to[0,\infty]$ given by $$C(x,\rho):=\begin{cases}
        0 \quad\;\; \delta_x \ld\rho,\\
        \infty \quad \text{otherwise,}
    \end{cases}$$
to get
$$\inf_{\pi\in\Pi(\mu,\nu)} \int \chi_{\delta_x\ld\pi_x}d\mu(x)= \sup_{\substack{\psi \in E \\ }} \mu(\psi) - \nu(\psi).$$
\end{proof}

Let us now recall the following relation for $\mu,\nu\in\P_2(\R^d)$:
$$\inf_{\eta\lc\nu} W^2_2(\mu,\eta)= \inf_{\pi\in\Pi(\mu,\nu)} \int_{\R^d} |x-\mean(\pi_x)|^2 d\mu(x),$$
where $W_2(\cdot,\cdot)$ is the 2-Wasserstein distance.
We prove an analogous equality for general cost functions on compact Polish spaces and relate it to our dual representation result.

\begin{theorem} \label{proj} Assume that $Y$ is compact and $E$ is a semi-stable cone of functions on $Y$. Let $C:X\times\PP(Y)\to[0,\infty]$ be lsc and convex in the second argument. Then for $\mu\in\PP(X)$, $\nu\in \PP(Y)$, we have
    $$\inf_{\eta\ld\nu}  \inf_{\pi\in\Pi(\mu,\eta) }\int C(x,\pi_x)d\mu(x)=\inf_{\pi\in\Pi(\mu,\nu)} \int \inf_{\xi\ld\pi_x}C(x,\xi) d\mu(x)=\sup_{\psi\in E}\mu(\psi^C)-\nu(\psi).$$
\end{theorem}
\begin{proof}
  Let us prove the first equality.  Fix $\varepsilon>0,$ and let $\eta^*\ld\nu$ and $\pi^*\in\Pi(\mu,\eta^*)$ be almost optimizers:
    \begin{align*}
        \inf_{\eta\ld\nu} \left[ \inf_{\pi\in\Pi(\mu,\eta) }\int C(x,\pi_x)d\mu(x)\right]+\varepsilon\geq
          \int C(x,\pi^*_x)d\mu(x).
    \end{align*}
By Theorem \ref{dil}, we have that $\eta^*\ld\nu$ implies the existence of a coupling $\widetilde{\pi}\in \Pi(\eta^*,\nu)$ that satisfies $\delta_{\widetilde{y}}\ld\widetilde{\pi}_{\widetilde{y}}$ for $\eta^*$-a.e. $\widetilde{y}\in Y$.

Define the probability measure $\pi$ on $X\times Y$ by $d\pi(x,y):=d\widetilde{\pi}_{\widetilde{y}}(y) d\pi^*(x,\widetilde{y}).$ One directly verifies that $\pi\in\Pi(\mu,\nu)$ and satisfies $\pi_x^*\ld\pi_x$ for $\mu$-a.e. $x\in X$. Hence
 \begin{align*}
          \int C(x,\pi^*_x)d\mu(x)\geq \int \inf_{\xi\ld\pi_x}C(x,\xi) d\mu(x) \geq \inf_{\pi\in\Pi(\mu,\nu)}\int\inf_{\xi\ld\pi_x}C(x,\xi) d\mu(x) .
    \end{align*}


For the converse inequality, fix $\varepsilon>0$ and let $\pi^*\in\Pi(\mu,\nu)$ be almost optimal and $\{\xi_x^*\}_{x\in X}$ a measurable selection of almost optimizers satisfying $\xi_x^*\ld \pi_x^*$. Then
$$\inf_{\pi\in\Pi(\mu,\nu)} \int \left[\inf_{\xi\ld\pi_x}C(x,\xi)\right] d\mu(x)+\varepsilon \geq \int C(x,\xi_x^*)d\mu(x).$$
Now define $d\eta(y):=\int_X d\xi_x^*(y)d\mu(x)$ and $d\pi(x,y):=d\xi^*_x(y)d\mu(x).$ Then $\eta\ld\nu$ and $\pi\in\Pi(\mu,\eta).$ This proves the first equality. 

To show the second equality, note that $\widehat{C}(x,\rho):=\inf_{\xi\ld\rho} C(x,\xi)$ is lsc (see \cite[Proposition 7.33]{bertsekas1978stochastic}), convex in $\rho$ and $\ld$-decreasing, so by Theorem \ref{mainthmgen}, we get
$$\inf_{\pi\in\Pi(\mu,\nu)} \int \inf_{\xi\ld\pi_x}C(x,\xi) d\mu(x)=\sup_{\psi\in E}\mu(\psi^{\widehat{C}})-\nu(\psi).$$
One easily shows that $\psi^{\widehat{C}}=\psi^C$ for $\psi\in E$, which finalizes the proof.
\end{proof}


    To illustrate the underlying idea of Theorem \ref{proj}, consider the following two-step process of transporting $\mu$ to $\nu$: first transport $\mu$ optimally (w.r.t. the cost function $C$)  to an intermediate measure $\eta$ that satisfies $\eta\ld\nu$, and then transport $\eta$ to $\nu$ via an $E$-dilation coupling, whose existence is guaranteed by Theorem \ref{dil}, with the intermediate measure $\eta$ chosen to minimize the transport cost between $\mu$ and $\eta$. The cost incurred by the first step (that is, the double infimum in Theorem \ref{proj}) is identical to the weak transport cost between $\mu$ and $\nu$ with respect to the $\ld$-decreasing hull of $C$, i.e. $\widehat{C}(x,\rho)=\inf_{\xi\ld\rho} C(x,\xi)$. Additionally, both values coincide with the dual formulation restricted to potentials in the semi-stable cone $E$.

    In the event that $C$ itself is $\ld$-decreasing, the cost function $\widehat{C}$ simplifies to $C$, the optimal $\eta$ is the target measure $\nu$, and the statement of Theorem \ref{proj} collapses to Theorem \ref{mainthmgen} for compact Polish spaces.






% \subsubsection{\color{black} Projections onto General Cones}
% \color{black}

% \begin{theorem}
%     Let $\mu\in\P(X),\, \nu\in\PP_p(Y)$ and let $C:X\times \P_p(Y) \to(-\infty,\infty]$ be jointly lsc, bounded from below and convex in the second argument. Then we have 
%     $$\inf_{\eta\ld\nu} \left[ \inf_{\pi\in\Pi(\mu,\eta) }\int C(x,\pi_x)d\mu(x)\right]=\inf_{\pi\in\Pi(\mu,\nu)} \int \left[\inf_{\xi\ld\pi_x}C(x,\xi)\right] d\mu(x).$$
% \end{theorem}
% \begin{proof}
%     Fix $\varepsilon>0.$ Let $\eta^*\ld\nu$ and $\pi^*\in\Pi(\mu,\eta^*)$ be almost optimizers:
%     \begin{align*}
%         \inf_{\eta\ld\nu} \left[ \inf_{\pi\in\Pi(\mu,\eta) }\int C(x,\pi_x)d\mu(x)\right]+\varepsilon\geq
%           \int C(x,\pi^*_x)d\mu(x).
%     \end{align*}
% By the generalized Strassen result, Theorem \ref{genStr}, we have that $\eta^*\ld\nu$ implies the existence of a coupling $\widetilde{\pi}\in \Pi(\eta^*,\nu)$ that satisfies $\delta_{\widetilde{y}}\ld\widetilde{\pi}_{\widetilde{y}}$ for $\eta^*$-a.e. $\widetilde{y}\in Y$.

% Define the probability measure $\pi$ on $X\times Y$ by $d\pi(x,y):=d\widetilde{\pi}_{\widetilde{y}}(y) d\pi^*(x,\widetilde{y}).$ One directly verifies that $\pi\in\Pi(\mu,\nu)$ and satisfies $\pi_x^*\ld\pi_x$ for $\mu$-a.e. $x\in X$. Hence
%  \begin{align*}
%           \int C(x,\pi^*_x)d\mu(x)\geq \int \inf_{\xi\ld\pi_x}C(x,\xi) d\mu(x) \geq \inf_{\pi\in\Pi(\mu,\nu)}\int\inf_{\xi\ld\pi_x}C(x,\xi) d\mu(x) .
%     \end{align*}


% For the converse inequality, fix $\varepsilon>0$ and let $\pi^*\in\Pi(\mu,\nu)$ be almost optimal, and $\{\xi_x^*\}_{x\in X}$ a measurable selection of almost optimizers satisfying $\xi_x^*\ld \pi_x^*$. Then
% $$\inf_{\pi\in\Pi(\mu,\nu)} \int \left[\inf_{\xi\ld\pi_x}C(x,\xi)\right] d\mu(x)+\varepsilon \geq \int C(x,\xi_x^*)d\mu(x)$$
% Now define $d\eta(y):=\int_X d\xi_x^*(y)d\mu(x)$ and $d\pi(x,y):=d\xi^*_x(y)d\mu(x).$
% \end{proof}



% \color{black}
% \fbox{\begin{minipage}{\textwidth}
%     Some potential cones to consider would be those corresponding to increasing functions, submodular functions, subharmonic functions, and plurisubharmonic functions (related to \cite{MR1035999}) amongst others. A study of costs of the form $C(x,\rho)=\int c(x,y)d\rho(y)$ is expected to recover the classical restriction to $c$-concave functions.
% \end{minipage}}


% \begin{itemize}
%     \item $\max$-stability of $D$: Notice that we only require $D$ to be a convex cone, containing the constants and $D\subseteq{\mathcal{C}}_p(Y)$. In particular, there is no assumption of $\max$-stability (i.e., $f,g\in D $ implies $f\lor g\in D$), despite this being a feature of $D$ in all three applications above.

%     One reason why one would potentially like to add the $\max$-stability assumption, is in hopes of proving a duality of the form 
% \begin{align*}
% 		\inf_{\pi\in\Pi(\mu,\nu)} \int C(x,\pi_x)d\mu(x)= \sup_{\psi \in {D}\cap( {\mathcal{C}}^{}_{b,p}+f_0)} \mu(R_C\psi) - \nu(\psi),
% 	\end{align*}
%  i.e., without $\bar D \cap({\mathcal{C}}^{lsc}_{b,p}+f_0)$,    which would sidestep the continuity argument for convex functions that we use in all three examples. I hoped that by taking a $\psi\in\bar D$ and approximating it by an increasing net $\{\psi_\alpha\}_{\alpha\in A}\subseteq D$ (in the spirit of \cite{ciosmak2024localisationconstrainedtransportsi} and \cite{aliprantis2007infinite}), I could pass to functions in $D$, but my approach required $\lim_{\alpha\in A}R_C (\psi_\alpha)=R_C(\psi)$, which I do not think holds. Perhaps a different approach would work.

%  Notice that, unlike the compact case I discussed in the previous draft, $D$ is cannot be assumed $\sup$-stable, which would imply $\bar D = D$, trivializing the above discussion. For example, in the setting of Subsection \ref{barcos}, take $f_a(y):=a|y|$ for $y\in\R^d$, $a\geq0$ then $f_a\in D$, but $\sup_{a\in(0,\infty]} f_a =  \infty {\chi}_{\R^d\setminus\{0\}}$, which is not real valued, hence cannot be in $D$.

%  On the other hand, perhaps this formulation is sufficient, since it allows for a general $D$, and then further refinements can be made on a case-by-case basis.
% \end{itemize}


% \appendix



% \section{Borel Measurability of \texorpdfstring{$A$}{A} in the Proof of Lemma \ref{rc stabil lemma}}
% \label{appa}
% \begin{proof}
%     Recall that $A$ is defined as
%     $$A:=\left\{(y,\xi)\in \R^d\times\PP_p(\R^d) \mid \conv{\psi_N}(y)+\varepsilon/2\geq \xi(\psi_N), \;\delta_y \lc\xi\right\}.$$
% To check Borel measurability, note that $A=F^{-1}([0,\infty))\cap G^{-1}(\{0\}),$ where $$F:\R^d\times\PP_p(\R^d)\to\R, \;\;G:\R^d\times\PP_p(\R^d)\to(-\infty,\infty]$$ are given by
% $$F(y,\xi):=\conv{\psi_N}(y)+\varepsilon/2-\xi(\psi_N),\;\;G(y,\xi):=\sup_{f\text{ convex}}f(y)-\xi(f).$$
% It is straightforward to show that $F$ is lower semicontinuous. On the other hand, $G$ is given by a supremum of continuous functions, so it is itself lower semicontinuous.
% \end{proof}









\bibliographystyle{plain}

\bibliography{joint_biblio, literatur}

%\bibliography{literatur}





\end{document}
